\clearpage\section{Lagrangeformalismen og bevarelseslove}\label{sec:lagrange}
\textbf{Fra opstilling til bevægelse.} I kapitel~1 valgte vi uafhængige koordinater, beregnede $T,V$ og opstillede $\mathcal L(q,\dot q,t)$. Hvordan bestemmes nu $q_i(t)$? Vi begynder med Hamiltons princip, udleder Euler--Lagrange og bruger derefter symmetrier, energi og bindinger til at reducere arbejdet.

\subsection{Virkningen og Hamiltons princip}
Til en mulig bane $q_i(t)$ mellem to faste tidspunkter $t_1,t_2$ knyttes virkningen
\[
S[q]=\int_{t_1}^{t_2}\mathcal L(q,\dot q,t)\,\dd t.
\]
Hamiltons princip siger, at den fysiske bane gør virkningen \emph{stationær}: $\delta S=0$ for alle tilstrækkeligt små variationer med faste endepunkter,
$\delta q_i(t_1)=\delta q_i(t_2)=0$. Stationær betyder ikke nødvendigvis et minimum.

Skriv $q_i(t,\varepsilon)=q_i(t)+\varepsilon\eta_i(t)$, hvor $\eta_i(t_1)=\eta_i(t_2)=0$. Den første variation er
\[
\delta S=\int_{t_1}^{t_2}\sum_i\left(
\frac{\partial\mathcal L}{\partial q_i}\delta q_i+
\frac{\partial\mathcal L}{\partial\dot q_i}\delta\dot q_i
\right)\dd t.
\]
Da $\delta\dot q_i=\dd(\delta q_i)/\dd t$, kan det andet led integreres ved dele:
\begin{align*}
\int_{t_1}^{t_2}\frac{\partial\mathcal L}{\partial\dot q_i}
\frac{\dd(\delta q_i)}{\dd t}\,\dd t
&=\left[\frac{\partial\mathcal L}{\partial\dot q_i}\delta q_i\right]_{t_1}^{t_2}
-\int_{t_1}^{t_2}\frac{\dd}{\dd t}
\left(\frac{\partial\mathcal L}{\partial\dot q_i}\right)\delta q_i\,\dd t\\
&=-\int_{t_1}^{t_2}\frac{\dd}{\dd t}
\left(\frac{\partial\mathcal L}{\partial\dot q_i}\right)\delta q_i\,\dd t,
\end{align*}
fordi randleddet er nul. Dermed
\[
\delta S=\int_{t_1}^{t_2}\sum_i\left[
\frac{\partial\mathcal L}{\partial q_i}-
\frac{\dd}{\dd t}\left(\frac{\partial\mathcal L}{\partial\dot q_i}\right)
\right]\delta q_i\,\dd t.
\]
Variationerne er vilkårlige indvendigt i tidsintervallet, så hver koefficient må forsvinde. Vi har således udledt Euler--Lagrange-ligningerne:
\key{\frac{\dd}{\dd t}\left(\frac{\partial\mathcal L}{\partial\dot q_i}\right)
-\frac{\partial\mathcal L}{\partial q_i}=0.}
Udledningen forudsætter tilstrækkeligt glat $\mathcal L$ og uafhængige variationer af de valgte generaliserede koordinater. En alternativ håndtering af bindinger er multiplikatormetoden i afsnit~\ref{sec:lambda}.

\subsection{Euler--Lagrange som fremgangsmåde}
Definér den generaliserede impuls $p_i=\partial\mathcal L/\partial\dot q_i$. Den lokale ligning er
$\dd p_i/\dd t=\partial\mathcal L/\partial q_i$; dette er ikke en påstand om, at $p_i$ generelt er bevaret.
For pendulet fra kapitel~1 har vi
\[
\mathcal L=\frac12ml^2\dot\theta^2+mgl\cos\theta,\quad
\frac{\partial\mathcal L}{\partial\dot\theta}=ml^2\dot\theta,\quad
\frac{\partial\mathcal L}{\partial\theta}=-mgl\sin\theta.
\]
Derfor
\[
ml^2\ddot\theta+mgl\sin\theta=0
\quad\Longrightarrow\quad
\boxed{\ddot\theta+\frac gl\sin\theta=0.}
\]
Dette er en \emph{differentialligning} for $\theta(t)$, ikke selve tidsløsningen. Hvis en ikke-konservativ generaliseret kraft $\mathcal Q_i^{\mathrm{(nk)}}$ ikke allerede er beskrevet i $\mathcal L$, gælder i stedet
\[
\boxed{\frac{\dd}{\dd t}\frac{\partial\mathcal L}{\partial\dot q_i}
-\frac{\partial\mathcal L}{\partial q_i}=\mathcal Q_i^{\mathrm{(nk)}}.}
\]
Eksempelvis giver lineær dæmpning $\mathcal Q^{\mathrm{(nk)}}=-b\dot q$ med $b>0$ et ekstra led i ligningen; man må ikke behandle almindelig friktion som et konservativt $V(q)$.
\begin{method}Vælg $q_i$ og $\mathcal L$, udregn først $\partial\mathcal L/\partial\dot q_i$, differentier derefter \emph{totalt} efter tiden, og træk $\partial\mathcal L/\partial q_i$ fra. Indsæt eventuelle ikke-konservative generaliserede kræfter, kontroller dimensioner/fortegn, og anvend initialdata til løsningen.\end{method}

\subsection{Den harmoniske oscillator i én og to dimensioner}
For et kvadratisk potentiale i én dimension,
\[
\mathcal L=\frac12m\dot q^2-\frac12kq^2,\qquad m,k>0,
\]
giver Euler--Lagrange $m\ddot q+kq=0$. Sæt $\omega=\sqrt{k/m}$; ligningen har løsningen
\key{q(t)=A\cos(\omega t)+B\sin(\omega t).}
Givet $q(0)=q_0$ og $\dot q(0)=v_0$ bliver $A=q_0$ og $B=v_0/\omega$. En generel begyndelsestid $t_0$ håndteres tilsvarende ved at skrive løsningen med $(t-t_0)$.

For en isotrop oscillator i et plan er
\[
\mathcal L=\frac12m(\dot x^2+\dot y^2)-\frac12m\omega^2(x^2+y^2),
\quad \ddot x+\omega^2x=0,\ \ddot y+\omega^2y=0.
\]
Begge komponenter har samme frekvens:
\[
x=A\cos\omega t+B\sin\omega t,\qquad
y=C\cos\omega t+D\sin\omega t.
\]
Som vektorligning er $\vv r(t)=\vv a\cos\omega t+\vv b\sin\omega t$. Derfor er banen billedet af en enhedscirkel under den lineære afbildning med kolonnerne $\vv a$ og $\vv b$: i almindelighed en ellipse med centrum i origo. Er $\vv a$ og $\vv b$ lineært afhængige, degenererer ellipsen til et linjestykke; er de ortogonale med samme længde, bliver banen en cirkel.
\begin{center}
\begin{tikzpicture}[scale=.80]
\draw[->,gray!70](-3.2,0)--(3.3,0)node[right]{$x$};
\draw[->,gray!70](0,-2.0)--(0,2.0)node[above]{$y$};
\draw[accent,very thick,rotate=25] (0,0) ellipse (2.45 and 1.17);
\fill (0,0) circle (1.3pt) node[below right,scale=.82]{$O$};
\draw[->,thick,red!70!black] (0,0)--(1.8,.72) node[above right]{$\vv r(t)$};
\node[accent,scale=.84] at (-2.7,1.55){Ellipse};
\end{tikzpicture}
\end{center}
\textbf{Fortolkning.} Det er parameteren $t$, der gennemløber ellipsen; den særlige oscillator--Kepler-forbindelse senere i pensum anvender generelt \emph{ikke} den samme tidsparameter. Banens form alene afgør derfor ikke tidsudviklingen.
\begin{method}Genkend et kvadratisk potentiale, aflæs frekvensen, løs for hver uafhængig koordinat og fastlæg alle konstanter fra startposition og starthastighed. Ved koblede oscillatorer kan normalmodemetoder være nødvendige.\end{method}

\subsection{Symmetrier, cykliske koordinater og impulsmoment}\label{sec:angular}
\textbf{Kan vi finde bevarede størrelser uden at løse ligningerne?} Hvis $\mathcal L$ er uafhængig af en generaliseret koordinat $q_i$, siger Euler--Lagrange direkte
\[
\frac{\partial\mathcal L}{\partial q_i}=0
\quad\Longrightarrow\quad
\frac{\dd}{\dd t}\frac{\partial\mathcal L}{\partial\dot q_i}=0
\quad\Longrightarrow\quad \boxed{p_i=\mathrm{konstant}}.
\]
En sådan $q_i$ kaldes \emph{cyklisk}. Sammenhængen mellem symmetri og bevarelse er bredere: et kontinuerligt transformationsinvariant virkningsprincip giver Noether-bevarelseslove også i koordinater, hvor symmetrien ikke fremstår som en enkelt cyklisk variabel.
\begin{center}\small\renewcommand{\arraystretch}{1.17}
\begin{tabular}{p{40mm}p{49mm}p{51mm}}\toprule
Symmetri & Simpel lokal test & Bevaret størrelse\\\midrule
Translation i $x$ & $\partial\mathcal L/\partial x=0$ & $p_x=\partial\mathcal L/\partial\dot x$\\
Rotation om $z$ & $\partial\mathcal L/\partial\phi=0$ & $p_\phi=\partial\mathcal L/\partial\dot\phi$\\
Tidsforskydning & $\partial\mathcal L/\partial t=0$ & $h=\sum_i\dot q_i p_i-\mathcal L$\\\bottomrule
\end{tabular}\end{center}
I kartesiske inertielle koordinater med mekanisk impuls $\vv p=m\dot{\vv r}$ er impulsmomentet omkring origo
\key{\vv L=\vv r\times\vv p,\qquad
L_x=yp_z-zp_y,\quad L_y=zp_x-xp_z,\quad L_z=xp_y-yp_x.}
\begin{center}
\begin{tikzpicture}[scale=.78,>=stealth]
\draw[->,gray](-.3,0)--(4.2,0)node[right]{$x$};
\draw[->,gray](0,-.3)--(0,2.6)node[above]{$y$};
\draw[->,thick,accent](0,0)--(2.65,1.75)node[above right]{$\vv r$};
\draw[->,thick](2.65,1.75)--(3.63,2.3)node[right]{$\vv p$};
\draw[accent](1.15,0) arc (0:33.4:1.15);
\node[accent] at (1.39,.42){$\phi$};
\node at (1.1,-.65){$L_z=xp_y-yp_x$};
\end{tikzpicture}
\end{center}
Ved en central kraft er drejningsmomentet $\vv\tau=\vv r\times\vv F=0$, så $\dd\vv L/\dd t=0$. I polære koordinater og uden ekstra hastighedsafhængige led er $L_z=p_\phi=mr^2\dot\phi$ for én partikel; for tolegemer bruges den reducerede masse i kapitel~\ref{sec:central}. Skeln mellem kanonisk og mekanisk impuls, hvor de ikke er ens.
\begin{method}Find koordinater, som ikke forekommer eksplicit i $\mathcal L$; differentier efter deres hastigheder for at få $p_i$. Ved rotationssymmetri: fastlæg akse, kontroller drejningsmoment eller cyklisk vinkel, og angiv den bevarede komponent af $\vv L$. Test symmetrien frem for automatisk at antage alle tre komponenter bevaret.\end{method}

\subsection{Energifunktionen og tidsinvarians}
Definér
\[
h(q,\dot q,t)=\sum_i\dot q_i\frac{\partial\mathcal L}{\partial\dot q_i}-\mathcal L.
\]
Differentier totalt:
\begin{align*}
\frac{\dd h}{\dd t}
&=\sum_i\left(\ddot q_i\frac{\partial\mathcal L}{\partial\dot q_i}
+\dot q_i\frac{\dd}{\dd t}\frac{\partial\mathcal L}{\partial\dot q_i}\right)
-\sum_i\left(\frac{\partial\mathcal L}{\partial q_i}\dot q_i
+\frac{\partial\mathcal L}{\partial\dot q_i}\ddot q_i\right)
-\frac{\partial\mathcal L}{\partial t}\\
&=\sum_i\dot q_i\left(\frac{\dd}{\dd t}
\frac{\partial\mathcal L}{\partial\dot q_i}
-\frac{\partial\mathcal L}{\partial q_i}\right)
-\frac{\partial\mathcal L}{\partial t}.
\end{align*}
På løsninger af Euler--Lagrange med ingen ikke-konservative ydre generaliserede kræfter forsvinder parenteserne, og derfor
\key{\frac{\dd h}{\dd t}=-\frac{\partial\mathcal L}{\partial t}.}
En tidsuafhængig Lagrangefunktion bevarer altså $h$. Hvis $\mathcal L=T-V$ med $V$ uafhængig af hastighederne og $T$ homogen af grad to i $\dot q_i$ uden andre hastighedsled, giver Eulers homogenitetssætning $\sum_i\dot q_i(\partial T/\partial\dot q_i)=2T$ og dermed $h=T+V$. Det gælder ikke ukritisk i reonome koordinater, systemer med hastighedsafhængige potentialer eller ved eksplicit tidsafhængighed.
\begin{method}Beregn $h=\sum_i\dot q_ip_i-\mathcal L$, ikke blot $T+V$ per automatik. Kontrollér $\partial\mathcal L/\partial t$ og eventuelle ikke-konservative kræfter, brug $h=E$ med initialbetingelser og kontrollér energiens dimensionsenhed. Den senere Hamiltonfunktion opstår ved Legendretransformation, når impulserne kan inverteres for hastighederne (kapitel~5).\end{method}
Ved ét konservativt koordinat kan $E=\frac12m\dot q^2+V(q)$ integreres for at finde $q(t)$; fortegnsvalg og vendepunkter er udledt som \emph{energiintegration} i kapitel~\ref{sec:central}.

\subsection{Ækvivalente Lagrangefunktioner og total tidsafledt}
To Lagrangefunktioner af formen
\[
\mathcal L_2(q,\dot q,t)=\mathcal L_1(q,\dot q,t)+
\frac{\dd F(q,t)}{\dd t}
\]
giver samme Euler--Lagrange-ligninger med faste konfigurationsendepunkter, fordi
\[
S_2-S_1=\int_{t_1}^{t_2}\frac{\dd F}{\dd t}\dd t
=F(q(t_2),t_2)-F(q(t_1),t_1),
\]
og variationen af dette randled er nul.
For at \emph{afgøre} ækvivalens skal forskellen sammenlignes med
\key{\frac{\dd F}{\dd t}=
\sum_i\frac{\partial F}{\partial q_i}\dot q_i+
\frac{\partial F}{\partial t}.}
Fremgangsmåden er at beregne $\Delta\mathcal L=\mathcal L_2-\mathcal L_1$, indsamle koefficienterne foran de uafhængige hastigheder, integrere dem til en kandidat $F(q,t)$ med en vilkårlig funktion af de resterende koordinater og tiden, og så kontrollere \emph{alle} koefficienter samt det hastighedsuafhængige led. Lokalt kræver dette kompatible krydsafledte mellem hastighedskoefficienterne.

\textbf{Eksempel.} Givet
\[
\Delta\mathcal L=tV'(q)\dot q+V(q)
\]
har vi $\partial F/\partial q=tV'(q)$. Integrationen efter $q$ giver
$F(q,t)=tV(q)+f(t)$. Kravet $\partial F/\partial t=V(q)$ medfører $f'(t)=0$ og derfor
\[
\boxed{F(q,t)=tV(q)+C,\qquad
\frac{\dd F}{\dd t}=tV'(q)\dot q+V(q).}
\]
Hvis forskellen er en total tidsafledt, er det en \emph{tilstrækkelig} betingelse for samme bevægelsesligninger; hvis testen fejler, følger det ikke automatisk, at ligningerne er forskellige. Impulserne kan forskydes:
\[
p_i^{(2)}=p_i^{(1)}+\frac{\partial F}{\partial q_i}.
\]
Dette giver en naturlig forbindelse til de senere genererende funktioner i kapitel~6.
\begin{method}Beregn $\Delta\mathcal L$; sammenlign hastighedsled med $F_{q_i}$; integrér og medtag vilkårlige funktioner af øvrige variable; brug de resterende koefficienter og $F_t$ til at fastlægge dem; differentier hele kandidaten og kontrollér lighed med $\Delta\mathcal L$.\end{method}

\subsection{Bindinger og Lagrange-multiplikatorer}\label{sec:lambda}
Kapitel~1 viste, at holonome bindinger ofte kan elimineres med koordinatvalg. Hvis det er uhensigtsmæssigt, kan man bevare redundante koordinater og pålægge $f_\alpha(q,t)=0$ ved multiplikatorer $\lambda_\alpha$.
For den angivne fortegnskonvention
$\mathcal L'=\mathcal L+\sum_\alpha\lambda_\alpha f_\alpha$,
med variation også i $\lambda_\alpha$, fås for ideelle holonome bindinger
\key{\frac{\dd}{\dd t}\frac{\partial\mathcal L}{\partial\dot q_i}
-\frac{\partial\mathcal L}{\partial q_i}
=\sum_\alpha\lambda_\alpha\frac{\partial f_\alpha}{\partial q_i},
\qquad f_\alpha(q,t)=0.}
Multiplikatorerne giver den generaliserede bindingskraft; deres numeriske værdi afhænger af bindingernes normalisering. For pendulet er vinklen $\theta$ én fri koordinat, men man kan i stedet arbejde med $(x,y)$ og $x^2+y^2-l^2=0$ samt én $\lambda$. For kontaktproblemer, hvor kontakt kan ophøre, skal $N\ge0$ kontrolleres; ved $N=0$ må den pålagte binding eventuelt ophæves.
\begin{method}Opstil $\mathcal L$ i de valgte (eventuelt redundante) koordinater; tilføj $\sum_\alpha\lambda_\alpha f_\alpha$; skriv de modificerede Euler--Lagrange-ligninger og bindinger; differentier om nødvendigt bindingerne; løs for koordinater og multiplikatorer; oversæt til fysiske bindingskræfter og tjek kontaktbetingelser.\end{method}

\subsection{Metodevalg og overgang til bevægede referencesystemer}
\begin{center}\small\renewcommand{\arraystretch}{1.16}
\begin{tabular}{p{58mm}p{79mm}}\toprule
Opgavespørgsmål & Første metode\\\midrule
Find bevægelsesligningen & Euler--Lagrange; tilføj eventuelle $\mathcal Q_i^{\mathrm{(nk)}}$\\
Find bevarede impulser & Cykliske koordinater eller Noethers symmetri\\
Find energibevarelse & $h=\sum\dot q_ip_i-\mathcal L$ og $\partial\mathcal L/\partial t$\\
Vis to Lagrangefunktioner ækvivalente & Søg en total tidsafledt $\dd F/\dd t$\\
Find kræfter fra ikke-eliminerede bindinger & Lagrange-multiplikatorer\\\bottomrule
\end{tabular}\end{center}
Vi har indtil nu opstillet bevægelsen i et inertielt system eller i en allerede angivet koordinatbeskrivelse. Hvad ændres, når selve origo accelererer, eller koordinatbasen roterer? Det undersøges i kapitel~4 om ikke-inertielle referencesystemer; her må også tidsafledte af koordinatbasen medtages.
